\chapter{Conclusion}
\label{chap:conclusion}

\section{Synthèse des contributions}

Ce mémoire a développé un système de détection précoce du DT1 adapté aux populations camerounaises.
Contributions : (1) dataset synthétique réaliste, (2) comparaison 6 algorithmes ML, (3) validation biomarqueurs via SHAP, (4) application démonstrative.

\section{Réponse aux objectifs}

Objectifs atteints : [RÉSUMER vos résultats principaux].
Hypothèses validées/partiellement validées : [BILAN].

\section{Perspectives de recherche}

Court terme : validation données réelles, amélioration recall.
Moyen terme : intégration clinique, formation personnel soignant.
Long terme : déploiement national, extension autres pays africains.

\section{Impact potentiel}

Amélioration diagnostic précoce, réduction mortalité infantile DT1, contribution santé publique camerounaise.

\vspace{1cm}

\begin{center}
\textit{``La science n'a pas de patrie, mais l'homme de science doit penser d'abord aux besoins de son pays."}\\
--- Louis Pasteur
\end{center}
